% !TeX program = xelatex
% UTF-8; standalone XeLaTeX/XePersian document, also usable with \input.
% Faithful transcription of Notebook/headache & migraine, including margins.
% Source photographs (unchanged; in chronological/page order):
% images/pic01.jpg = 20261016_204810057.jpg
% images/pic02.jpg = 20261016_204818262.jpg
% images/pic03.jpg = 20261016_204823130.jpg
% images/pic04.jpg = 20261016_204832527.jpg
% images/pic05.jpg = 20261016_204838147.jpg
% images/pic06.jpg = 20261016_204842325.jpg
% Medical wording and statements are transcribed, not clinically revised.
\ifdefined\settextfont
  \def\HMstandalone{0}
\else
  \def\HMstandalone{1}
  \documentclass[11pt,a4paper]{article}
  \usepackage[margin=2cm]{geometry}
  \usepackage{amsmath,amssymb,enumitem}
  \usepackage{xepersian}
  \IfFontExistsTF{Amiri}{\settextfont{Amiri}}{\settextfont{FreeSerif}}
  \setlatintextfont{FreeSerif}
  \setlist[itemize]{itemsep=3pt,topsep=4pt,parsep=0pt}
\fi
\providecommand{\en}[1]{\lr{#1}}
\newcommand{\HMpage}[1]{%
  \clearpage\section*{سردرد و میگرن — صفحهٔ #1}
  {\footnotesize\noindent\en{Subject: 180}\hfill\en{Date: \quad / \quad /}\par}}
\newcommand{\HMfooter}[1]{\par\medskip{\footnotesize\noindent\en{YASHA}\hfill #1\par}}
\ifnum\HMstandalone=1 \begin{document}\fi

\HMpage{۱}
\begin{latin}\noindent\textbf{Headache and Migraine}\end{latin}
\subsection*{اتیولوژی}
\begin{latin}
\noindent Headache
\begin{itemize}
  \item \textbf{primary}
  \begin{itemize}
    \item Migraines
    \item Tension-type headache
  \end{itemize}
  \item \textbf{Secondary}
  \begin{itemize}
    \item Associated with minor illnesses
    \begin{itemize}\item URI, Sinusitis\end{itemize}
    \item First symptoms of serious conditions
    \begin{itemize}\item Meningitis\item Brain tumors\end{itemize}
  \end{itemize}
\end{itemize}
\end{latin}
\subsection*{تظاهرات بالینی}
\noindent\en{Temporal pattern}: برای \en{DDX} سردرد:
\begin{latin}
\begin{itemize}
  \item Acute
  \item Recurrent episodic
  \item chronic progressive
  \item chronic non-progressive
\end{itemize}
\end{latin}
\subsection*{تنشن تایپ}
\begin{itemize}
  \item \textbf{شایع‌ترین نوع} سردرد راجعهٔ اولیه در کودکان و نوجوانان.
  \item عموماً خفیف و بدون علائم همراه.
  \item \textbf{فعالیت روزانه را مختل نمی‌کند.}
\end{itemize}
\begin{latin}
\begin{itemize}
  \item The pain is \textbf{global} and \textbf{squeezing} or \textbf{pressing}.
  \item \textbf{For hours or days}.
\end{itemize}
\end{latin}
\begin{itemize}
  \item \textbf{علائم همراه وجود ندارد}: تهوع، استفراغ، فتوفوبی، فونوفوبی.
  \item می‌تواند با استرس‌های تحصیلی، اضطراب، افسردگی مرتبط باشد.
  \item حاشیهٔ متصل به «اضطراب، افسردگی»: در زمینهٔ سردرد یک اضطراب یا افسردگی باشد.
\end{itemize}
\HMfooter{۱}

\HMpage{۲}
\begin{latin}
\noindent\textbf{Migraine headaches}
\begin{itemize}
  \item A type of \textbf{recurrent} headache
  \item Begin in childhood
  \item Pain:
  \begin{itemize}
    \item \textbf{Frontal} or \textbf{bitemporal} or \textbf{unilateral}
    \item Moderate to Severe
    \item Pounding or Throbbing
  \end{itemize}
\end{itemize}
\end{latin}
\begin{itemize}
  \item درد: \textbf{با فعالیت تشدید می‌شود}.
  \item \en{\textbf{last 1--72 hours}}
  \item \textbf{علائم همراه} در میگرن: تهوع، استفراغ، رنگ‌پریدگی، فتوفوبی، فونوفوبی.
  \item \textbf{عوامل تخفیف}: مکان تاریک، ساکت.
  \item میگرن در نوپایان: اپیزودهای تحریک‌پذیری، بی‌قراری و رنگ‌پریدگی؛ تهوع، استفراغ.
  \item \textbf{اورا} می‌تواند قبل و یا همراه با میگرن رخ بدهد.
  \item بین \textbf{۵ تا ۶۰ دقیقه} ادامه می‌یابد.
  \item \textbf{اوراهای بینایی} شایع‌ترین اوراها هستند.
  \item تعدادی اوراها کامپلکس، آتیپیکال یا سردرد: \textbf{نیاز به بررسی}؛ \en{EEG, Neuroimaging}.
  \item موارد متصل به اوراهای کامپلکس/آتیپیکال:
  \begin{latin}
  \begin{itemize}
    \item Hemiparesis, confusion,
    \item Vertigo, Monocular blindness
  \end{itemize}
  \end{latin}
\end{itemize}
\HMfooter{۲}

\HMpage{۳}
\begin{itemize}
  \item علل شایع سردرد ثانویه: تروما سر، بیماری ویروسی، سینوزیت.
  \item سردرد ناشی از مصرف بیش‌ازحد دارو: \en{Medication-Overuse}.
  \item می‌تواند سردردهای اولیه و ثانویه را کامپلیکه کند.
  \item علل جدی سردرد ثانویه:
  \begin{itemize}
    \item \en{$\uparrow$ICP} (تومور، مالفورماسیون عروقی).
    \item \textbf{سودوتومور سربری}.
  \end{itemize}
\end{itemize}
\subsection*{چه زمانی به \en{$\uparrow$ICP} شک می‌کنیم؟}
\begin{itemize}
  \item \textbf{سردرد و استفراغ} همراه آن بعد از بیدار شدن یا بعد از \textbf{بلند شدن از خواب}، یا \textbf{بعد از سرفه} و مانور والسالوا تشدید می‌شود، یا خم شدن به جلو.
  \item \textbf{ادم پاپیل و یا \en{FND}}: \textbf{بررسی جدی علت سردرد}.
\end{itemize}
\subsection*{اقدامات تشخیصی}
\begin{itemize}
  \item در اکثر کودکان با شرح حال و \en{Ph/E} به تشخیص می‌رسیم و نیاز به بررسی‌ها نیز نیست.
  \item \textbf{تصویربرداری معمولاً ضرورتی ندارد.}
\end{itemize}
\subsubsection*{در چه مواردی تصویربرداری انجام می‌دهیم؟ \en{MRI}}
\begin{itemize}
  \item معاینهٔ نورولوژیک غیرطبیعی.
  \item شک به \en{$\uparrow$ICP}.
  \item \textbf{اورا آتیپیک}.
  \item بدتر شدن سردرد در طول زمان و سیر بیماری. % Handwritten margin at bottom left.
\end{itemize}
\HMfooter{۳}

\HMpage{۴}
\subsection*{چه زمانی \en{CT} می‌کنیم؟}
\begin{itemize}
  \item \textbf{سردرد ناگهانی و بسیار شدید}: بررسی \en{ICH}.
  \item در این موارد اگر \en{CT} منفی بود: \textbf{\en{LP}}.
  \begin{itemize}
    \item بررسی فشار \en{Opening}.
    \item پلئوسیتوز، \en{$\uparrow$RBC}، زانتوکرومیا.
    \item بررسی \en{SAH}.
  \end{itemize}
\end{itemize}
\subsection*{درمان}
\begin{itemize}
  \item جزء اصلی درمان میگرن با اپیزودهای \en{infrequent}: \textbf{داروهای ضد درد به‌صورت \en{intermittent}}.
  \item تجویز زودهنگام آنالژزیک و استراحت در یک اتاق \textbf{ساکت، تاریک}.
  \item دارو:
  \begin{itemize}
    \item \textbf{استامینوفن}.
    \item \textbf{\en{NSAID}}: ناپروکسن، ایبوپروفن.
  \end{itemize}
  \item هیدریشن و داروهای ضد تهوع هم کمک‌کننده هستند.
  \item اگر خط اول کافی نبود: \en{Triptan agents}.
  \begin{itemize}
    \item \textbf{آگونیست گیرندهٔ سروتونین}.
    \item فرم: تزریقی، قرص، اسپری بینی.
    \item ریسک \en{Stroke} را بالا می‌برند.
  \end{itemize}
\end{itemize}
\HMfooter{۴}

\HMpage{۵}
\subsection*{کنتراندیکاسیون تریپتان‌ها}
\begin{itemize}
  \item میگرن همی‌پلژیک.
  \item میگرن با اورا \en{brainstem}.
\end{itemize}
\subsection*{قبل از شروع درمان دارویی در میگرن}
\begin{itemize}
  \item تغییر لایف استایل.
  \item مصرف بیشتر مایعات.
  \item پرهیز از نخوردن وعده‌های غذایی.
  \item بهداشت خواب.
  \item ورزش.
  \item پرهیز از مواد تحریک‌کننده.
\end{itemize}
\subsection*{مکمل‌های کمکی}
\begin{itemize}
  \item \textbf{ریبوفلاوین}.
  \item \textbf{منیزیم}.
\end{itemize}
\subsection*{اندیکاسیون درمان با \en{Daily Preventive Agents}}
\begin{itemize}
  \item کودکان با \textbf{بیش از یک حملهٔ ناتوان‌کننده در هفته}.
  \item ضدافسردگی سه‌حلقه‌ای: \en{Amitriptyline, Nortriptyline}.
  \item ضدتشنج‌ها: \en{Topiramate, Valproic acid}.
  \item آنتی‌هیستامین: \en{Cyproheptadine}.
  \item بتابلاکر: پرانول.
  \item \en{CCB}: \en{Verapamil}.
  \item داروی جدید: \en{Calcitonin gene-related peptide receptor antagonists}.
\end{itemize}
\HMfooter{۵}

\HMpage{۶}
\subsection*{چهار الگوی اصلی سردرد در کودکان}
\subsubsection*{\en{a) Acute}}
\begin{itemize}
  \item اولین اپیزود سردرد بدون سابقهٔ قبلی.
  \item اگر بسیار شدید: در بزرگسالان \textbf{\en{SAH}} ناشی از آنوریسم.
  \item در کودکان: دنبال تب در \en{URI}؛ [واژهٔ خط‌خورده: ویروسی].
\end{itemize}
\subsubsection*{\en{b) Acute recurrent}}
\begin{itemize}
  \item حمله‌های راجعه با یک \textbf{اینتروال بدون علامت}.
  \item الگوی شایع در سندرم‌های سردرد اولیه: \textbf{میگرن}، \textbf{\en{TTH}}.
  \item ممکن است در برخی سندرم‌های صرعی دیده شود:
  \en{\textbf{Benign occipital epilepsy}}.
\end{itemize}
\subsubsection*{\en{c) Chronic progressive}}
\begin{itemize}
  \item افزایش شدت و فرکانس سردرد در طول زمان.
  \item \textbf{شک به \en{$\uparrow$ICP}}: \textbf{سودوتومور سربری}، هیدروسفالی، تومور مغزی، آبسه، مننژیت.
\end{itemize}
\subsubsection*{\en{d) Chronic nonprogressive}}
\begin{itemize}
  \item \textbf{بیشتر از ۳ ماه} که در هر ماه \textbf{حداقل ۱۵ سردرد}.
  \item معاینهٔ نورولوژیک \en{NL} است.
  \item \textbf{عوامل روان‌پزشکی} شایع است.
\end{itemize}
\HMfooter{۶}

\ifnum\HMstandalone=1 \expandafter\enddocument\fi
